% LeafletFA Application Note: status for David (facts only; not manuscript text). Built 2026-11-15.
% Sources: OUTLINE_PREP.md (budgets, 2026-11-13), LEAFLETFA_SCREEN.md (word counts at 78436ab),
% ~/LeafletFA/scripts/validation/VALIDATION_PREREG.md and VALIDATION_REPORT.md (through ceb5048),
% figure AN = paper_figures/figAN/figAN.pdf (built 2026-11-13), scaling = LeafletFA_validation/scaling/.
\documentclass[10pt]{article}
\usepackage[a4paper,margin=16mm]{geometry}
\usepackage[T1]{fontenc}\usepackage{helvet}\renewcommand{\familydefault}{\sfdefault}
\usepackage{graphicx,booktabs,tabularx,longtable,array,xcolor,enumitem}
\usepackage[hidelinks]{hyperref}
\setlist{nosep,leftmargin=*}
\setlength{\parindent}{0pt}\setlength{\parskip}{4pt}
\newcolumntype{L}[1]{>{\raggedright\arraybackslash}p{#1}}
\newcommand{\m}{$-$}
\newcommand{\gee}{$\geq$\,}
\renewcommand{\arraystretch}{1.15}
\begin{document}

{\Large\bfseries LeafletFA Application Note: status}\\[2pt]
For David \quad\textbar\quad 15 November 2026 \quad\textbar\quad facts only, no manuscript text

\begin{enumerate}
  \item Current figure (figure AN)
  \item The Application Note, section by section: word budget and state
  \item Every validation check, 10.1 to 15b: pre-registered rule, result, verdict
  \item Scaling and runtime
\end{enumerate}

\textbf{Framing question.} After the simulation (15b), the note reads as an atlas resource and a fast estimator for sparse
data, not as a differential splicing method. Options on the table: (a) keep it a Bioinformatics Application Note on that
narrower claim; (b) fold LeafletFA into the SpliceVI paper as a baseline and data resource.
What the checks support and do not support is in section 3.

\section*{1. Figure AN (built 13 Nov 2026; 183 $\times$ 164 mm)}
\begin{center}\includegraphics[width=\textwidth]{/gpfs/commons/groups/knowles_lab/Karin/LeafletFA_bioinformatics/paper_figures/figAN/figAN.pdf}\end{center}
\begin{itemize}
  \item a: one catalog over two atlases (TMS 78,135 cells; Allen 54,571), event types; b: accuracy against held-out
    cells (check 10.1) with the check 14 shrinkage comparator; c--f: \textit{Gria2} flip/flop, \textit{Fgfr2} IIIb/IIIc,
    \textit{Nrxn1} SS4 (the true 90-nt SS4 exon), \textit{Ptprc}.
  \item Done this week: combined catalog column; complex events split (linked mutually exclusive 1,034; other 7,056);
    programs removed from the figure.
  \item Open: panel a model schematic (placeholder); which version of panel b; caption not written.
  \item Supplement: known-biology gallery built (10 rows; one caption fact stale). No figure or text yet for scaling,
    the interval and accuracy checks, or the simulation.
\end{itemize}

\clearpage
\section*{2. The Application Note, section by section}
Word counts: prose only, comments and red text excluded (screen of \texttt{main.tex} at 78436ab, 11 Nov; text since
then only wrapped stale lines, no rewording). Budget 2,000 words for abstract + body.

\footnotesize
\begin{longtable}{L{36mm}rrL{98mm}}
\toprule
Section & Budget & Now & State \\
\midrule
Abstract & 150 & 15 & Empty apart from the supplementary line. \\
Introduction & 200 & 256 & Written, about 55 over. One stale sentence on programs (programs failed checks 3 and 8). \\
Data and event catalog & 250 & 221 & Written; text shared word for word with SpliceVI. Alignment (STAR two-pass), junction calling (regtools) and filters are in the supplement. The filter-chain table, which SpliceVI cites, is empty. \\
ATSEmapper & 150 & 156 & Written. Inputs, outputs and run time missing. Its event definition is not yet reconciled with the shared catalog text that SpliceVI uses. \\
Model and fitting & 200 & 267 & Written, about 65 over. The E-step and the program embedding that SpliceVI uses as a baseline are not in the main text; the supplementary model section is a TODO. \\
Cell-type usage and intervals & 150 & 77 & Stale: describes the check 10 interval. The shipped interval is check 12. \\
Results: catalog & 120 & 42 & Facts in comments, almost no prose. \\
Results: accuracy (10.1) & 180 & 11 & Facts in comments, almost no prose. \\
Results: limits & 170 & 31 & Facts in comments, almost no prose. \\
Results: known biology, portal & 120 & 129 & Framing only. \textit{Gria2}, \textit{Fgfr2}, \textit{Nrxn1} SS4 and \textit{Ptprc} need sentences. Bebee 2015 reference missing. Portal has no text. \\
Results: scale & 60 & 77 & Numbers wrong. Production fits took about 13.2 min (TMS) and 38.5 min (Allen), not 10.6 / 31.6 (those are projection times). CPU about 100$\times$ slower than GPU. GPU-memory figures unsourced. \\
Discussion & 250 & 230 & Mostly written. Needs the opening and the SpliceVI companion sentence. \\
\midrule
\textbf{Total} & \textbf{2,000} & \textbf{about 1,520} & Filling the under-budget Results sections to budget adds about 600 words, which gives about 2,100. Cuts of about 190 words of duplicated text are already identified. \\
\midrule
Availability, data availability & -- & -- & Outside the budget. Need URLs, the deposit description and the Zenodo DOI. \\
\bottomrule
\end{longtable}
\normalsize

\textbf{What the note has to carry, and where it sits now.}
\begin{itemize}
  \item \textbf{Data processing} (alignment, junction calling, filters, cells kept): main text = shared catalog paragraph;
    details in the supplement; filter-chain table empty.
  \item \textbf{ATSE mapping} (ATSEmapper): main-text section written; inputs, outputs and run time missing.
  \item \textbf{Model} (Dirichlet-multinomial factor model, $K = 20$, SVI recipe, interval): main text written;
    interval paragraph stale; full model, E-step and embedding not yet written in the supplement.
  \item \textbf{Application} (catalog, accuracy, known biology, portal, scale): figure AN built; prose mostly missing.
\end{itemize}

\textbf{Points where SpliceVI depends on this paper.}
\begin{itemize}
  \item SpliceVI cites LeafletFA for the processing pipeline and for the filter-chain table (empty here).
  \item SpliceVI describes LeafletFA as a baseline (E-step against frozen programs, centered-log embedding, live-program
    rule); none of that is in the LeafletFA main text.
  \item The two papers' ATSE definitions need one wording (SpliceVI: shared donor/acceptor component with a 1\% filter and
    mutually exclusive exon split; LeafletFA ATSEmapper paragraph: splice-graph component with event types).
  \item LeafletFA does not yet mention SpliceVI in its main text.
\end{itemize}

\textbf{Not part of the paper, needed for release:} browser URL (not published), tutorial plan, deposit.

\clearpage
\section*{3. Validation checks 10.1 to 15b}
Rules were written and committed before each run; verdicts are as recorded. There is no check 11. Simulation depths:
full / 30\% / 10\% of real reads. ``Raw'' = TRAIN pseudobulk; ``B'' = shrinkage of each cell type to the all-cells
$\psi$ with one tuned constant (check 14). TV = total variation error.

\scriptsize
\begin{longtable}{L{23mm}L{42mm}L{70mm}L{26mm}}
\toprule
Check & Pre-registered rule & Result (TMS / Allen, or full / 30\% / 10\%) & Verdict \\
\midrule
\endhead
10.1 Cell-type $\psi$ vs held-out pseudobulk & LeafletFA mean TV $\leq$ raw, and paired CI $<0$ in the smallest-type bin or the 1--19-read bin (amend.\ 6) & TMS 0.0940 vs 0.1012, $\Delta$ \m0.0072 [\m0.0094, \m0.0050]; 1--19 reads \m0.0625. Allen 0.0625 vs 0.0647, $\Delta$ \m0.0021 [\m0.0039, \m0.0005]; 1--19 reads \m0.0543 & \textbf{PASS / PASS} \\
10.2 Cell-type-specific calls replicate across mouse halves & Confirmed fraction \gee0.5 and \gee5$\times$ the permuted-label null & Confirmed 0.813 / 0.837; null 0.422 / 0.771 (1.9$\times$ / 1.09$\times$) & \textbf{FAIL / FAIL} \\
10.2 LeafCutter control & Descriptive; cannot change 10.2 & Same design, LeafCutter vs LeafletFA: real 0.521 vs 0.819 (TMS), 0.740 vs 0.858 (Allen); null 0.348 vs 0.422, 0.604 vs 0.795 & Neither method meets the rule \\
10.2c Call sets across seeds & Median Jaccard \gee0.7 (reported) & 0.637 / 0.770 & fail / pass \\
10.3 Interval coverage (check 10 interval) & Reported against nominal 0.95 & 0.14--0.35 across read bins and atlases & Far below nominal; replaced by check 12 \\
10.4 \textit{Fgfr2} IIIb & Every epithelial type above every mesenchymal type, and IIIb specific in \gee1 epithelial type & Epithelial min 0.704 $>$ mesenchymal max 0.649; specific in 10 epithelial types & \textbf{PASS} \\
10.4 \textit{Ptprc} & \gee1 B/T call with direction confirmed in TEST & 27 calls, 25 confirmed & \textbf{PASS} \\
10.4 \textit{Nrxn1--3} SS4 & \gee1 neuronal pair with TV $>\delta$ and non-overlapping intervals & 1,256 / 946 / 1,129 of 1,891 pairs (true SS4, note 6) & \textbf{PASS} $\times$3 \\
10.5 Catalog export & Specification only & 6,478,290 / 4,420,578 rows & done \\
10.6 Runtime, atlas size & Descriptive & 132,706 cells: fit 73.2 min, job 119.4 min, 40.7 GB host, 25.0 GB GPU & descriptive \\
10.7 Pairwise cell-type comparisons & As 10.2, pooled over type pairs & TMS, 69 pairs: confirmed 0.780, null 0.473 (1.65$\times$) & \textbf{FAIL} \\
12 Predictive + recalibrated interval & Coverage in [0.90, 1.00] in both read bins, both atlases, both directions (amend.\ 7) & All eight cells 0.950--0.963 & \textbf{PASS} (shipped) \\
13 Group-pooled estimate & Adopt only if top-100 error is lower in both atlases and the overall CI is not $>0$ in both & Top-100: 0.253 vs 0.224 / 0.100 vs 0.123; overall TMS +0.0064 [+0.0042, +0.0084] & Not adopted \\
14 Shrinkage comparator B & LeafletFA $-$ B CI $<0$ in both atlases (amend.\ 9) & TMS \m0.0029 [\m0.0044, \m0.0014]; Allen +0.0006 [\m0.0008, +0.0019] & \textbf{Not met} (TMS better, Allen tie) \\
15 Simulation, accuracy (amend.\ 10) & Descriptive; primary = cell types with their own truth & $\Delta$ raw +0.0118 / +0.0140 / +0.0140 & Worse than raw and B \\
15 Simulation, differential & Descriptive & AUROC 0.975 / 0.970 / 0.966; LeafCutter on shared units 0.962 / 0.971 / 0.982 & Inflated: one program separated the halves (AUROC 1.00), a design flaw \\
15b Half check (amend.\ 11) & Best single-program AUROC: $\leq$0.60 none, 0.60--0.75 partial, $>$0.75 confounded & 0.711 / 0.664 / 0.692 & Partial \\
15b Accuracy, own-truth types & $\Delta$ CI rule; compared with check 15 & $\Delta$ raw +0.0097 / +0.0128 / +0.0125; $\Delta$ B +0.0105 / +0.0147 / +0.0163; gap ratio 0.82--0.91 of check 15 & \textbf{Worse} than raw and B; gap survives in part \\
15b Accuracy, all pairs & Same rule (secondary) & $\Delta$ raw \m0.0060 / \m0.0060 / \m0.0106; at 1--19 reads \m0.0817 (full); $\Delta$ B +0.0085 / +0.0129 / +0.0130 & Better than raw; worse than B \\
15b Differential, type-specific shifts & Descriptive, read with half-check caveat & Pooled AUROC 0.684 / 0.664 / 0.648 (raw 0.632 / 0.620 / 0.604) & Descriptive \\
15b vs LeafCutter (shared units) & CI of $\Delta$ excludes 0 & Per-type $\Delta$ \m0.027 / \m0.024 / \m0.011; LeafCutter covers 14 / 9 / 5\% of units & \textbf{LeafCutter better} at every depth \\
Scaling (15 Nov) & Descriptive & Section 4 & descriptive \\
\bottomrule
\end{longtable}
\normalsize

\textbf{Held up:} accuracy against held-out cells on real data (10.1, both atlases), with the largest gain where reads are thin
(1--19 reads); calibrated predictive intervals (12); the known-biology checks (10.4); speed and scaling.\\
\textbf{Did not hold up:} cell-type-specific calls against the permuted null (10.2, 10.7); improvement over simple
shrinkage in both atlases (14); accuracy where reads are plentiful, and differential ranking against LeafCutter, in
simulation (15b).

\clearpage
\section*{4. Scaling and runtime}
\textbf{Existing records.}
\footnotesize
\begin{longtable}{L{34mm}L{15mm}L{20mm}L{30mm}L{28mm}L{24mm}}
\toprule
Run & Job & Cells & Hardware & Time & Peak memory \\
\midrule
LeafletFA TMS fit (16,200 steps) & 22400619 & 54,694 & RTX PRO 6000 Blackwell, 8 cores & 13.2 min; 0.05 s/step & 9.2 GB host; GPU not recorded \\
LeafletFA Allen fit (16,200 steps) & 22400622 & 38,199 (of 54,571) & same & 38.5 min; 0.14 s/step & 21.3 GB host; GPU not recorded \\
LeafletFA atlas-size fit & 22597359 & 132,706 & same & 73.2 min fit; 0.271 s/step & 40.7 GB host; 25.0 GB GPU \\
scQuint & 22725257 & 54,694 & 16 CPUs & 86.0 min (20 epochs) & 1.7 GB \\
LeafCutter catalog benchmark & 22780123 & per-type pseudobulk & 128 tasks $\times$ 4 CPUs & 19.2 task-h; 77 core-h & 1.28 GB per task \\
\bottomrule
\end{longtable}

\textbf{Controlled runs} (TMS TRAIN subsets, 2,000 steps, minibatch min(8,192, cells); all GPU runs on one RTX PRO 6000 Blackwell).
\begin{longtable}{lllrrrr}
\toprule
Method & Device & Cells & s/step & Time (min) & Peak GPU (GB) & Peak host (GB) \\
\midrule
LeafletFA & GPU & 5,000 & 0.032 & 1.28 & 2.51 & 2.07 \\
LeafletFA & GPU & 10,000 & 0.046 & 1.62 & 4.14 & 2.64 \\
LeafletFA & GPU & 25,000 & 0.046 & 1.63 & 5.11 & 4.82 \\
LeafletFA & GPU & 50,000 & 0.046 & 1.63 & 6.72 & 8.73 \\
LeafletFA & GPU & 54,694 & 0.046 & 1.62 & 7.02 & 8.92 \\
LeafletFA & CPU (16) & 5,000 & 3.27 & 110.1 & -- & 3.65 \\
scQuint & CPU (16) & 5,000 & -- & 13.3 (60 epochs) & -- & 1.17 \\
scQuint & CPU (16) & 10,000 & -- & 37.3 (55) & -- & 1.18 \\
scQuint & CPU (16) & 25,000 & -- & 72.3 (40) & -- & 1.38 \\
scQuint & CPU (16) & 50,000 & -- & 67.0 (25) & -- & 1.77 \\
\bottomrule
\end{longtable}
\normalsize
\begin{itemize}
  \item CPU is 102$\times$ slower than GPU per step at 5,000 cells. Time per step is flat from 10,000 cells up (fixed minibatch).
  \item LeafletFA times are for a fixed 2,000 steps; scQuint times run to its own early stop. They are not the same quantity.
\end{itemize}
\begin{center}\includegraphics[width=\textwidth]{/gpfs/commons/groups/knowles_lab/Karin/LeafletFA_validation/scaling/scaling.pdf}\end{center}

\end{document}
